\section{Design of LLM-GridEval}
\label{sec:design}

\subsection{Framework Architecture}
The LLM-GridEval framework integrates the system model into a unified evaluation environment through three layers (Fig.~\ref{fig:framework-architecture}).

\begin{figure*}
    \centering
    \includegraphics[width=0.5\linewidth]{figures/system.png}
    \caption{High level architecture of the LLM-GridEval framework, showing the co-simulation, sandboxed interface, and LLM-based cognitive layers.}
    \label{fig:framework-architecture}
\end{figure*}

\textbf{Co-simulation Layer}: Implements $\text{Sim}(\cdot)$ through HELICS federates with synchronized time advancement: a GridPACK transmission federate (IEEE 9-bus), two GridLAB-D distribution federates (IEEE 123-node feeders; Feeder A contains EV loads, Feeder B contains only background loads), a feeder controller federate implementing threshold-based defense, and an MCP federate connecting the EV management interface.

\textbf{Interface and Tool Layer}: Acts as a sandboxed API boundary between co-simulation and attacker. Each attack primitive is defined by a schema for inputs/outputs, preconditions, and mapping to HELICS operations. The layer performs argument validation and constraint enforcement (e.g., EV capacity bounds, rate limits) before invoking HELICS calls, ensuring all actions are physically realizable and auditable.

\textbf{Cognitive Layer}: Hosts the attacker policy $\pi$ using a ReAct-like interaction loop: (1) observe grid state via \texttt{get\_grid\_status()}, (2) construct prompt with state and available primitives, (3) query LLM for action specification, (4) validate against schema and constraints, (5) execute via sandboxed interface, (6) advance simulation. This closed loop connects LLM reasoning to physically meaningful grid actions while the constrained catalog mitigates hallucinated commands.

This layered separation serves two purposes. First, it enables controlled experimentation: researchers can swap attacker implementations (scripted, RL-based, or LLM-based) while keeping the co-simulation and tool interface fixed, isolating the impact of policy changes. Second, it provides security guarantees: the sandboxed interface ensures that even a fully autonomous LLM agent cannot issue physically invalid commands or bypass rate limits, making automated red-teaming safe for deployment against production-representative grid models.

\subsection{Prototype Implementation}
The prototype instantiates LLM-GridEval for EV capacity attacks on a transmission-distribution co-simulation.

\textbf{Physical Plant}: The IEEE 9-bus transmission system connects to two IEEE 123-node distribution feeders via HELICS value federates with conservative time advancement \cite{hardy2024helics}. Feeder A hosts six EV charging stations (EV1--EV6) distributed across phases A, B, and C at different nodes \cite{ieee123feeder}. EV1 and EV4 are co-located with local storage and can be islanded from the main feeder. Feeder B contains only uncontrollable background loads for demand diversity.

\textbf{MCP Server}: Implemented as a HELICS federate, the MCP server exposes the two attack primitives. \texttt{GET /status} implements \texttt{get\_grid\_status()} and returns feeder power at the swing node, individual EV capacities and actual powers, and islanding states. \texttt{POST /set\_capacity} implements \texttt{set\_ev\_capacity(ev\_id, P)}, validates requests against configured capacity bounds, and publishes setpoints to GridLAB-D. This design mirrors realistic EV fleet management systems where cloud-hosted services provide programmatic control interfaces.

\textbf{Feeder Controller}: The defender is a rule-based controller that samples total feeder power $P_{\text{feeder}}(t)$ at the feeder-head (swing node). When $P_{\text{feeder}} \geq P_{\max}$ (4.2~MW), the controller sheds EV load; when $P_{\text{feeder}} \leq P_{\min}$ (2.6~MW), it restores EV capacities toward nominal levels. This threshold-based control provides a baseline defense for evaluation.

\textbf{LLM Orchestration}: At configurable simulation intervals, a coordination module queries the MCP server via \texttt{get\_grid\_status()} and formats the response into a structured prompt. The prompt includes the current grid state, available primitives with their schemas, attack constraints, and high-level objectives. The LLM returns a JSON action specification that is validated against the primitive schema before execution. For the experiments in this paper, we use a local LLM endpoint (Llama-3.1-8B) to ensure reproducibility and avoid API rate limits during extended campaigns.
